\section{发表与获奖情况}

本作品此前未在期刊、会议、预印本平台或其他公开出版物发表，未登记为既有获奖成果，也未以当前完整作品形态参加其他比赛。本次为首次提交。与作品配套的技术文档、作品介绍、演示视频和界面截图均为本次赛事制作，没有将其他项目的材料改名后重复使用。

代码仓库应本次比赛要求公开，地址为 \textbf{\repo}。仓库为 Public，保留按功能模块和核心 AI 会话划分的提交记录。历史用于核对实现顺序和返工过程，不为参赛补造提交，不修改时间戳。仓库中的 README、技术文档、测试、评测结果和 AI 对话快照可与本文互相核对。

作品当前没有论文、专利、软件著作权或奖项可列示。本文不以“已申请”“在审”或“拟发表”等表述填补空白。若后续产生正式发表或获奖记录，应以可核验的公开链接、证书编号或官方页面另行更新，本节目前只陈述截至 2026 年 11 月 18 日的事实。

\begin{table}[H]
\centering\footnotesize
\caption{公开材料与核验入口}
\begin{tabularx}{\textwidth}{@{}L{3.0cm}Y L{3.5cm}@{}}
\toprule
\textbf{材料} & \textbf{位置} & \textbf{可核对内容} \\
\midrule
公开源码 & \repo & 业务实现、细粒度历史、许可与提交说明 \\
复现指南 & 仓库根目录 \file{README.md} & 环境要求、启动命令、配置、功能与边界 \\
技术文档 & \file{docs/技术文档.md} 及配套 PDF & 架构、意图控制、接口、测试、评测和局限 \\
AI 对话快照 & \file{docs/ai-logs/} & 关键提示词、人工采纳与否决、落地结果 \\
测试与评测 & \file{apps/server/tests/}、\file{apps/web/tests/}、\file{scripts/eval-results/} & 确定性断言、逐题或逐问句结果、统计口径 \\
\bottomrule
\end{tabularx}
\end{table}
